\chapter*{Апстракт}
\addcontentsline{toc}{chapter}{Апстракт}

Процесот на лично логирање често е фрагментиран низ повеќе апликации и бара рачно внесување на податоци во однапред дефинирани форми.
Во овој дипломски труд се развива Logche, локално-прв систем што ги претвора кратките неструктурирани записи во структурирани податоци погодни за пребарување, филтрирање и анализа.
Истражувањето ги опфаќа подготовката и чистењето на податочни множества, fine-tuning и benchmark постапка за споредба на основни и дообучени модели.
Qwen моделот се приспособува за шест категории: храна, повеќекомпонентни оброци, тренинг, финансии, медиуми и движење.
Посебно внимание се посветува на локалната обработка, приватноста, конзистентниот JSON излез и поделбата на одговорностите меѓу јазичниот модел и детерминистичкиот апликациски код.
